\section*{第 \ref{ch:g6:inside-the-earth} 章——地球内部}

\begin{solution}{exo:g6:inside-the-earth:1}
\omterm{def:g6:inside-the-earth:layers}{地壳}、\omterm{def:g6:inside-the-earth:layers}{地幔}、\omterm{def:g6:inside-the-earth:layers}{外核}、\omterm{def:g6:inside-the-earth:layers}{内核}。\omterm{def:g6:inside-the-earth:layers}{外核}是
液态的。\omterm{def:g6:inside-the-earth:layers}{外核}和\omterm{def:g6:inside-the-earth:layers}{内核}由金属（主要是铁，还有一些
镍）组成；\omterm{def:g6:inside-the-earth:layers}{地壳}和\omterm{def:g6:inside-the-earth:layers}{地幔}由\omterm{def:g3:stones-pebbles-sand:rock}{岩石}组成。
\end{solution}

\begin{solution}{exo:g6:inside-the-earth:2}
（a）错误：\omterm{def:g3:continents-and-oceans:continent}{大陆}之下的\omterm{def:g6:inside-the-earth:layers}{地壳}（约 30~km）比
\omterm{def:g2:water-on-earth:ocean}{海洋}之下（约 5~km）的\omterm{def:g6:inside-the-earth:layers}{地壳}要厚。（b）错误：\omterm{def:g6:inside-the-earth:layers}{地核}由金属构成，
是铁和镍的合金。（c）正确。（d）正确：12\,262~m，约 12.3~km。
\end{solution}

\begin{solution}{exo:g6:inside-the-earth:3}
（a）2891~km，约 2900~km。（b）5150~km。（c）$5150 - 2891 = 2259$：
约 2260~km。
\end{solution}

\begin{solution}{exo:g6:inside-the-earth:4}
“蛋壳像\underline{\omterm{def:g6:inside-the-earth:layers}{地壳}}，蛋白像
\underline{\omterm{def:g6:inside-the-earth:layers}{地幔}}，蛋黄像\underline{\omterm{def:g6:inside-the-earth:layers}{地核}}。但与蛋黄不同，
\underline{\omterm{def:g6:inside-the-earth:layers}{地核}}分为两部分：液态的\underline{\omterm{def:g6:inside-the-earth:layers}{外核}}
和固态的\underline{\omterm{def:g6:inside-the-earth:layers}{内核}}。”
\end{solution}

\begin{solution}{exo:g6:inside-the-earth:5}
$8.849 \div 6371 \approx 0.00139$，且 $0.00139 \times 100 \approx 0.14$：
珠穆朗玛峰仅约为地球半径的 \textbf{0.14~\%}。从远处看，
地球是一个非常光滑的球体：它最高的山峰只是一个微小的凸起。
\end{solution}

\begin{solution}{exo:g6:inside-the-earth:6}
（a）$30 \div 5 = 6$：厚了六倍。（b）$100 \div 5 = 20$：厚了
二十倍。
\end{solution}

\begin{solution}{exo:g6:inside-the-earth:7}
（a）$12 + 1 \times 25 = 37~^\circ$C。（b）$12 + 1.5 \times 25 = 12 + 37.5 =
49.5~^\circ$C。（c）$12 + 2 \times 25 = 62~^\circ$C。
\end{solution}

\begin{solution}{exo:g6:inside-the-earth:8}
温度升高 $74 - 14 = 60$ 度，历时 2~km，故 $60 \div 2 =
\textbf{30}$~度每千米。是的：30 介于东部平均值
（25）与西部平均值（34）之间。
\end{solution}

\begin{solution}{exo:g6:inside-the-earth:9}
（a）刚好在\omterm{def:g6:inside-the-earth:layers}{地核}之上，约为每秒 13.7~km；刚好在其顶部之下，
约为每秒 8.1~km：P 波突然减速。（b）在
2891~km 与 5150~km 之间，即\omterm{def:g6:inside-the-earth:layers}{外核}中。（c）S 波不能在
液体中传播：\omterm{def:g6:inside-the-earth:layers}{外核}是液态的。
\end{solution}

\begin{solution}{exo:g6:inside-the-earth:10}
（a）$12.3 \div 30 = 0.41$：约为平均
\omterm{def:g3:continents-and-oceans:continent}{大陆}\omterm{def:g6:inside-the-earth:layers}{地壳}的 \textbf{41~\%}。（b）不能：在\omterm{def:g3:continents-and-oceans:continent}{大陆}之下，\omterm{def:g6:inside-the-earth:layers}{地幔}始于\omterm{def:g6:inside-the-earth:layers}{地壳}之下，
平均深约 30~km；钻孔在\omterm{def:g6:inside-the-earth:layers}{地壳}中停止，远未到达\omterm{def:g6:inside-the-earth:layers}{地幔}。
\end{solution}

\begin{solution}{exo:g6:inside-the-earth:11}
（a）$29 \times 6371 = 184\,759$，且 $15 + 184\,759 = 184\,774$：约
185\,000~度。（b）实验室给出的\omterm{def:g6:inside-the-earth:layers}{内核}边缘温度约为 6000~度；
莱奥的数值约大了 30~倍（
$184\,774 \div 6000 \approx 31$）。（c）每千米 29~度的升温率
仅在最初约 2~km 内测得。莱奥将其用于数千
千米之外的测量范围；在深处，温度升高要
慢得多。速率只能用于测量过的范围之内。
\end{solution}

\begin{solution}{exo:g6:inside-the-earth:12}
（a）若\omterm{def:g6:inside-the-earth:layers}{地壳}比平均值轻，则深处的物质必须
比 5.5 更密，才能把平均值拉高。（b）位于\omterm{def:g6:inside-the-earth:layers}{外核}和
\omterm{def:g6:inside-the-earth:layers}{内核}中（密度约从 9.9 到 13~克每立方厘米）；
\omterm{def:g6:inside-the-earth:layers}{地幔}在其最底部才刚好达到 5.6。（c）密度较大的物质会沉到
较轻物质之下，正如石头在水中下沉：重金属最终位于
中心，较轻的\omterm{def:g3:stones-pebbles-sand:rock}{岩石}位于上方。
\end{solution}

\begin{solution}{pb:g6:inside-the-earth:1}
\textbf{1.} $2 \times 6371 = 12\,742$~km。模型宽 2~m，即
2000~mm。

\textbf{2.} $12\,742 \div 2000 \approx 6.37$：模型上 1~mm 代表
真实地球约 6.37~km。

\textbf{3.} $8.849 \div 6.371 \approx 1.39$：在模型上，珠穆朗玛峰约
高 \textbf{1.4~mm}，相当于一粒\omterm{def:g3:stones-pebbles-sand:sand}{沙子}。

\textbf{4.} $12.3 \div 6.371 \approx 1.9$~mm。模型中心位于
表面以下 1000~mm（1~m）处：钻孔约
短 520 倍（$1000 \div 1.9 \approx 520$）。

\textbf{5.} $2891 \div 6.371 \approx 453.8$：\omterm{def:g6:inside-the-earth:layers}{外核}顶部在
表面以下约 \textbf{454~mm} 处涂色。

\textbf{6.} $5150 \div 6.371 \approx 808.3$：约在表面以下 808~mm
处。\omterm{def:g6:inside-the-earth:layers}{内核}在模型上的半径为 $1000 - 808 = 192$~mm。

\textbf{7.} \omterm{def:g6:inside-the-earth:layers}{地幔}：从约 5~mm 到 454~mm，约 449~mm。\omterm{def:g6:inside-the-earth:layers}{外
核}占 $808 - 454 = 354$~mm，\omterm{def:g6:inside-the-earth:layers}{内核}占 192~mm。

\textbf{8.} $30 \div 6.371 \approx 4.7$~mm：一条细带，可以
用细毛笔涂色。

\textbf{9.} $15 + 2 \times 29 = 73~^\circ$C。在模型上，2~km 相当于
表面以下 $2 \div 6.371 \approx 0.31$~mm：几乎只有
三张纸的厚度。

\textbf{10.} $4 \div 6.371 \approx 0.63$~mm，而钻孔约为
1.9~mm：矿井较浅，约浅 $1.9 - 0.6 = 1.3$~mm。

\textbf{11.} 物质是否熔化，取决于温度和
压强。\omterm{def:g6:inside-the-earth:layers}{内核}所受的挤压最大，约为
空气压强的 330 万倍：受压的铁只有在
更高的温度下才会熔化，因此\omterm{def:g6:inside-the-earth:layers}{内核}中的铁虽比\omterm{def:g6:inside-the-earth:layers}{外核}更热，
却仍保持固态。

\textbf{12.} $5 \div 6.371 \approx 0.785$：在 2~m 宽的模型上，大洋\omterm{def:g6:inside-the-earth:layers}{地壳}
只需涂成 \textbf{0.78~mm} 厚，相当于
约八张纸的厚度（$0.78 \div 0.1 \approx 8$）。
\end{solution}
